% Avsnitt 19.10, ur lectures/L19/appendix/c_exercises.md.

\section{Övningar}
\label{c19:sec:exercises}

\begin{aside}
Övning~\ref{c19:ex:thirdnode}, \ref{c19:ex:feature} och \ref{c19:ex:afterabort} befäster
\secref{c19:sec:verify}. Övning~\ref{c19:ex:bridge} till \ref{c19:ex:endtoend} befäster
\secref{c19:sec:bridge} till \ref{c19:sec:gaps} och bilaga~\ref{app:protocol}, som är auktoriteten
överallt där den och en övning är oense.
\end{aside}

\exercise{En tredje nod, och ett nytt arbitreringsutfall}{Simulering}
\label{c19:ex:thirdnode}
\xpart{a} Utöka \file{can_controller_tb.vhd} med en tredje nod, C, som delar samma \code{bus_line}.
Ge C ID:t \code{0x002} och låt den begära sändning på samma cykel som nod A (\code{0x000}) och B
(\code{0x001}). Bekräfta (genom simulering, inte bara förutsägelse) att A fortfarande vinner, och att
\emph{både} B och C upptäcker arbitreringsförlust.

\xpart{b} Nod A:s \code{0x000} är den lägsta identifierare en standardram kan bära, så inget val av
ID låter C slå den. Höj A:s ID till \code{0x010}, och välj sedan ett ID till C som gör C till vinnare
medan B (\code{0x001}) fortfarande förlorar. Motivera ditt val med \secref{c10:sec:arbitration}:s
arbitreringsregel, och bekräfta sedan den nya vinnaren i simulering. En följd att lägga märke till
innan dina assertions överraskar dig: med A på \code{0x010} slår B:s \code{0x001} även A, så i den
här omgången \textbf{förlorar A också} och höjer \code{error}; bara C fullföljer.

\exercise{En funktion ur en riktig kontroller, som designskiss}{Resonemang}
\label{c19:ex:feature}
Välj \textbf{en} post ur \secref{c19:sec:gaps}:s lista över vad riktiga CAN-kontrollrar gör som den
här designen inte gör. Skissa, utan att nödvändigtvis implementera den i VHDL (i prosa, plus ett
tillägg i tillståndsdiagrammet eller en ny portlista om det hjälper), vad som skulle behöva ändras i
\code{can_controller} för att stödja den. Ange vilket eller vilka befintliga block
(\code{bit_timer}, \code{crc15}, \code{tx_shift_reg}/\code{rx_shift_reg}, eller
\code{can_controller} självt) som skulle behöva ändras, och vilka som kunde stå exakt som de är.

\exercise{\code{spi_reg_bridge}}{VHDL}
\label{c19:ex:bridge}
Skriv din egen \file{spi_reg_bridge.vhd} utifrån \secref{c19:sec:bridge} och
bilaga~\ref{app:protocol}. Verifiera den:

\begin{shell}
cd bridge
ghdl -a --std=93 spi_def.vhd spi_reg_bridge.vhd spi_reg_bridge_tb.vhd
ghdl -e --std=93 spi_reg_bridge_tb
ghdl -r --std=93 spi_reg_bridge_tb --assert-level=error
\end{shell}

\xpart{a} Få alla fyra fallen att passera. Fall 3 och 4 är de två som gör den här modulen till mer än
en byteräknare \textemdash{} läs deras kommentarer innan du börjar, inte efteråt.

\xpart{b} Gör nu sönder den med flit. Ta bort varje användning av \code{ss_active} och räkna bytes
modulo fem i stället. Fall 1 och 2 passerar fortfarande. Säg vilken kontroll i fall 3 som fallerar,
citera värdet som felmeddelandet rapporterar, och förklara var de bytesen kom ifrån.

\xpart{c} En brygga som ignorerar de reserverade kommandobitarna passerar också fall 1 och 2.
Förklara varför protokollspecifikationen lägger regeln om reserverade bitar i \emph{det här} lagret i
stället för i \code{register_bank}, givet att banken aldrig ser mer än ett 4-bitars index.

\xpart{d} Läsvägen låser \code{reg_rdata} \textbf{en gång}, vid slutet av kommandobyten. Säg vad
\code{register_bank} skulle behöva göra i stället om bryggan läste det på nytt för varje databyte, och
varför det vore fel plats att skjuta ner regeln i banken. Ställ tillbaka modulen.

\exercise{Komponera \code{can_spi_node}}{VHDL}
\label{c19:ex:node}
Skriv toppnivån utifrån \secref{c19:sec:node}: \code{spi_slave}, \code{spi_reg_bridge},
\code{register_bank} och \code{can_controller}, plus en \code{meta_prev} för reseten.

\xpart{a} Skriv den, och kör \code{can_spi_node_tb} mot den. Fyra utdelade testbänkar täcker redan
allt \emph{inuti} filen; säg vilka fyra, och säg sedan den enda sortens fel som
\code{can_spi_node_tb} kan fälla och som ingen av de fyra kan. (Ledtråden är den positionella
portbindningen.)

\xpart{b} \code{can_spi_node_tb} håller 200~ns tyst kring varje \code{SS}-flank, där
bilaga~\ref{app:protocol} kräver minst 60~ns. Den är alltså snäll mot tidskraven med flit, och ingen
utdelad testbänk pressar dem. Kopiera den, krymp den tysta perioden stegvis ner mot 60~ns, och ta
reda på vid vilket värde din nod slutar svara.

Säg sedan \textbf{varför} just där, i klockcykler räknat, och var i \code{spi_slave} gränsen sitter.
Jämför det uppmätta värdet med de 60~ns specifikationen lovar: har du marginal, eller ligger din nod
precis på gränsen? Det här är enda stället i boken där du \emph{mäter} en marginal i stället för att
läsa den, och det är värt att göra noga.

\xpart{c} Räkna \code{meta_prev}-instanserna i den färdiga designen och säg vad var och en
synkroniserar. Säg sedan varför \code{can_spi_node} räcker \code{can_controller} den \emph{råa}
\code{reset_n} i stället för den \code{reset_s2_n} den ger varje annat block, och vad som skulle gå
fel om den synkroniserade den två gånger.

\xpart{d} \code{tx_bus} och \code{bus_en} lämnar designen som två skilda portar i stället för en
ledning. Förklara vad som måste hända mellan dem och en fysisk ledning som delas med ett andra kort,
och vad som skulle fallera om du tilldelade \code{tx_bus} rakt till en pinne och band ihop två kort.

\exercise{Vad kortet visade}{Resonemang}
\label{c19:ex:board}
Bring-upen kördes i åtta steg, i två halvor: tre på två DE0-CV och fem med hatten som master. Den här
övningen handlar om vad de fastställde och vad de inte gjorde. Den kräver ingen egen hårdvara.

\xpart{a} Stegen är ordnade så att varje steg utesluter en felklass innan nästa beror på den. Säg för
vart och ett av steg 1, 4 och 5 (ensamt kort, loopback på hatten, registereko) vilken enda sak det
bevisar som steget före inte gjorde. Säg sedan vilka av de åtta stegen som \code{can_spi_node_tb}
\textbf{redan} hade fastställt i simulering, och vilka som bara hårdvara kan fastställa.

\xpart{b} Två kopplingsfel att skilja åt. Säg vilket steg som fallerar först om MISO och MOSI är
förväxlade, och vilket som fallerar först om \code{VDDIO2} aldrig kopplades in, så att \code{PORTC}
står utan matning. Förklara varför de fallerar vid olika steg, och vad den skillnaden låter dig sluta
dig till utan att röra en mätare. Säg också vad \code{MVIO.STATUS} hade svarat i det andra fallet,
och varför det svaret är snabbare än mätaren.

\xpart{c} Testprogrammet som delades ut är skrivet för att lyckas: det håller \code{SS} låg i god
tid, skickar aldrig en sjätte byte och avbryter aldrig mitt i en transaktion. Säg vilken av
protokollspecifikationens regler \textemdash{} de 60~ns kring \code{SS}, \code{MISO} under
kommandobyten, eller avbrottsregeln \textemdash{} som \textbf{inget} av de åtta stegen skulle ha
avslöjat som trasig, och varför just den därför är den dyraste att ha fel i när en drivrutin du inte
skrivit möter noden.

\xpart{d} En ensam nod kan aldrig få \code{rx_valid} att pulsa, hur dess pinnar än är kopplade,
vilket är skälet till att passet lade en ram över \textbf{två} kort. Säg varför en nod aldrig kan ta
emot sin egen ram (skälet är \chapref{c17:ch}:s \code{role}), och vad det andra kortet därmed bevisar
som varken en loopback eller \code{can_controller_tb} kan.

\xpart{e} De två korten går på två oberoende 50 MHz-oscillatorer. \Secref{c19:sub:gaps-tb} kallar den
delade klockan i \code{can_controller_tb} för testbänkens \textbf{största} gap av precis det skälet.
Var noga med vad som ändras: \code{bit_timer.resync} gör redan riktigt arbete i simulering, så säg vad
de två oberoende klockorna lägger till som ingen simulering med delad klocka kan visa, och ungefär
hur långt isär de två oscillatorerna skulle kunna driva innan en ram slutade komma fram hel
(Övning~\ref{c12:ex:drift} gjorde den uträkningen; återanvänd den). Jämför den gränsen med en typisk
kristalls tolerans och säg om demonstrationen egentligen sätter den på prov.

\xpart{f} \code{rx_bus} kommer utifrån FPGA:n, och \code{can_spi_node} skickar den osynkroniserad
vidare till \code{can_controller} med flit, eftersom \code{can_controller} innehåller sin egen
\code{meta_prev} (\chapref{c12:ch}). Förklara vad som skulle kunna gå fel om den instansen togs bort,
varför felet skulle vara betydligt svårare att diagnostisera än en felkopplad ledning, och varför
ingen mängd simulering i den här boken någonsin skulle återskapa det.

\exercise{Att följa ett \code{send()} från ände till ände}{Resonemang}
\label{c19:ex:endtoend}
\xpart{a} Följ en enskild ramförfrågan hela vägen ner. Utgå från de fem bytesen på MOSI som skriver
\code{0x1} till \code{TX_SEND}, och följ dem genom \code{spi_slave} (bytes),
\code{spi_reg_bridge} (en registerskrivning) och \code{register_bank} (ett \code{tx_req} på en cykel)
fram till \code{can_controller}s portar. Namnge vad varje lager räcker vidare till nästa. Gör sedan
detsamma uppåt, för hur drivrutinen får veta att ramen är klar.

\xpart{b} Fyra lager är tre fler än \code{can_controller_tb} behöver. Namnge för varje gräns den
utdelade testbänk som spikar fast den, och säg vad som skulle vara oprövat om den testbänken inte
fanns.

\xpart{c} Pollning är det enda den här designen erbjuder: det finns ingen avbrottsutgång någonstans.
Säg vad som skulle behöva läggas till för att ge drivrutinen en sådan, ur vilken befintlig signal den
skulle härledas, och i vilket lager den skulle behöva ligga. Förklara sedan varför det hade varit
hopplöst för en drivrutin att polla ett \code{tx_done} som varar en cykel, och vilken modul som gjorde
att det inte är det.

\xpart{d} En drivrutin skriver \code{TX_ID}, sedan \code{TX_DLC}, sedan \code{TX_SEND}, som tre
separata transaktioner. Mellan den andra och den tredje går bussen till tomgång och en annan nod
sänder. Säg om något av det drivrutinen skrivit är i fara, och vilken regel i registerkartan det är
som gör ditt svar sant.

\exercise{Livet efter ett avbrott}{Simulering}
\label{c19:ex:afterabort}
\Secref{c19:sub:gaps-tb} listar ”ingen nod sänder någonsin igen efter ett avbrott” som något
systemtestbänken inte täcker. Den här övningen täcker det, och visar vad \code{clear}-ingången hos
\code{crc15} var till för.

\xpart{a} Utöka \file{can_controller_tb.vhd} med ett tredje fall. Ge båda noderna tid att återgå till
tomgång efter att fall 2 körts, alltså efter att nod B förlorat arbitreringen och övergett sin ram
mitt i identifieraren, och låt sedan \textbf{B} sända en egen ram (ID \code{0x123}, DLC 2, valfri
nyttolast) med A som mottagare. Bekräfta i simulering att A tar emot den: A:s \code{rx_valid} pulsar,
A:s \code{error} förblir låg, och ID:t kommer tillbaka korrekt.

\xpart{b} Ta nu bort nollställningen: ändra \code{crc_clear} till en konstant \code{'0'} och kör om.
Fall 1 och 2 passerar fortfarande. Fall 3 gör det inte. Säg vilken nod som höjer \code{error}, vid
vilket tillstånd, och förklara varför det \emph{inte} är den noden som har buggen.

\xpart{c} Räkna ut vad B:s \code{crc15}-register håller i det ögonblick den avbryter i fall 2, och
varför ingen mängd korrekt beteende därefter någonsin återför det till noll. (Ledtråden är
\chapref{c13:ch}:s egenskap att generering och kontroll är samma operation: vad förutsätter den om
registret vid en rams \emph{början}?)

\xpart{d} Hur många ramar måste nod B skicka innan den fungerar igen, med nollställningen fortfarande
borttagen? Förklara ditt svar, och säg vad en drivrutin som pollar \code{STATUS} och
\code{ERROR_FLAGS} över SPI skulle observera varje gång, med \chapref{c18:ch}:s regel för hur felbiten
låses.

\xpart{e} \Chapref{c16:ch} lägger nollställningen i \code{STATE_START} i stället för i
\code{STATE_CRC_WAIT} eller på vägen ut ur ett avbrott. Ge ett skäl vardera till varför de andra två
placeringarna är sämre.
